%% ============================================================
%% Part V — Proof-Carrying Python  (Slides 121 – 140)
%% ============================================================

%% --- Slide 121: Section divider ---
\begin{frame}
\frametitle{}
\vfill
\begin{center}
{\Large\textcolor{jugeoBlue}{\textbf{Part V}}}\\[6pt]
{\LARGE\textcolor{jugeoDark}{\textbf{Proof-Carrying Python}}}\\[10pt]
\begin{tikzpicture}
  \draw[jugeoBlue, line width=1.2pt] (0,0) -- (4,0);
\end{tikzpicture}\\[10pt]
{\large\textit{\textcolor{gray}{The proof ships with the code}}}
\end{center}
\vfill
\end{frame}

%% --- Slide 122: Why Not Extraction? ---
\begin{frame}
\frametitle{Why Not Extraction?}
\begin{itemize}
  \item F* extracts to \jugmark{OCaml}
  \item Lean extracts to \jugmark{C}
  \item Coq extracts to \jugmark{OCaml / Haskell}
\end{itemize}
\vspace{8pt}
\begin{alertblock}{The Problem}
  The shipped code carries \textcolor{jugeoRed}{\textbf{no proof metadata}}.\\[4pt]
  \emph{``How do you know this is correct?''}\\
  $\longrightarrow$ You point to a \textbf{different codebase} in a \textbf{different language}.
\end{alertblock}
\end{frame}

%% --- Slide 123: JuGeo's Approach — Skip Extraction ---
\begin{frame}
\frametitle{JuGeo's Approach --- Skip Extraction}
\begin{exampleblock}{Key Idea}
  Generate \jugmark{proof-carrying Python} \textbf{directly}.
\end{exampleblock}
\vspace{10pt}
\begin{itemize}
  \item Standard Python code
  \item Annotated with \textcolor{jugeoBlue}{verification metadata}
  \item The proof \jugmark{travels} with the code
  \item No separate artifact, no different language
\end{itemize}
\vspace{8pt}
\begin{center}
\begin{tikzpicture}
  \node[fill=jugeoGreen!20, rounded corners, inner sep=8pt, font=\large]
    {Code + Proof = \textbf{one Python file}};
\end{tikzpicture}
\end{center}
\end{frame}

%% --- Slide 124: The Generation Pipeline ---
\begin{frame}
\frametitle{The Generation Pipeline}
\begin{center}
\begin{tikzpicture}[
  stage/.style={draw=jugeoBlue, fill=jugeoLight!40, rounded corners=4pt,
                minimum width=2.4cm, minimum height=0.8cm,
                font=\scriptsize\bfseries, text=jugeoDark, align=center},
  arr/.style={-{Stealth[length=5pt]}, jugeoBlue, thick}
]
  \node[stage] (s1) {Specification};
  \node[stage, right=0.35cm of s1] (s2) {Cover\\Design};
  \node[stage, right=0.35cm of s2] (s3) {Inhabitant\\Fleets};
  \node[stage, right=0.35cm of s3] (s4) {Construction\\Loop};
  \node[stage, below=0.6cm of s1, xshift=1.4cm] (s5) {Descent\\and Glue};
  \node[stage, right=0.35cm of s5] (s6) {Certificate\\Emit};
  \node[stage, right=0.35cm of s6, fill=jugeoGreen!25] (s7) {Proof-Carrying\\Python};
  \draw[arr] (s1) -- (s2);
  \draw[arr] (s2) -- (s3);
  \draw[arr] (s3) -- (s4);
  \draw[arr] (s4) |- ([yshift=-0.15cm]s4.south) -| (s5);
  \draw[arr] (s5) -- (s6);
  \draw[arr] (s6) -- (s7);
\end{tikzpicture}
\end{center}
\vspace{4pt}
Seven stages from specification to deployable, verified Python.
\end{frame}

%% --- Slide 125: Stage 1 — Cover Design ---
\begin{frame}
\frametitle{Stage 1 --- Cover Design}
\begin{block}{Goal}
  Decompose the verification goal into manageable \jugmark{patches}.
\end{block}
\vspace{6pt}
\begin{center}
\begin{tikzpicture}[
  patch/.style={draw=jugeoBlue, fill=jugeoLight!30, rounded corners,
                minimum width=1.8cm, minimum height=1.2cm, font=\footnotesize},
  arr/.style={-{Stealth[length=4pt]}, jugeoBlue!70, dashed}
]
  \node[patch] (p1) at (0,0) {Patch A};
  \node[patch] (p2) at (2.5,0) {Patch B};
  \node[patch] (p3) at (5,0) {Patch C};
  \node[patch] (p4) at (1.25,-1.6) {Patch D};
  \node[patch] (p5) at (3.75,-1.6) {Patch E};
  \draw[arr] (p1) -- (p2);
  \draw[arr] (p2) -- (p3);
  \draw[arr] (p1) -- (p4);
  \draw[arr] (p2) -- (p4);
  \draw[arr] (p2) -- (p5);
  \draw[arr] (p3) -- (p5);
\end{tikzpicture}
\end{center}
\vspace{4pt}
Like dividing a large map into \textbf{overlapping sections}\\
that are easier to verify individually.
\end{frame}

%% --- Slide 126: Stage 2 — Inhabitant Fleets ---
\begin{frame}
\frametitle{Stage 2 --- Inhabitant Fleets}
Generate \jugmark{candidate implementations} from multiple sources:
\vspace{6pt}
\begin{center}
\begin{tikzpicture}[
  src/.style={draw=jugeoBlue, fill=jugeoLight!30, rounded corners,
              minimum width=2.4cm, minimum height=0.7cm, font=\footnotesize},
  pool/.style={draw=jugeoGreen!80, fill=jugeoGreen!15, rounded corners,
               minimum width=3cm, minimum height=0.8cm, font=\footnotesize\bfseries},
  arr/.style={-{Stealth[length=5pt]}, jugeoBlue, thick}
]
  \node[src] (a) at (0, 1.2) {Solver synthesis};
  \node[src] (b) at (0, 0.0) {Copilot suggestions};
  \node[src] (c) at (0,-1.2) {Runtime witnesses};
  \node[src] (d) at (0,-2.4) {Human input};
  \node[pool] (p) at (5.5,-.6) {Candidate Pool};
  \draw[arr] (a.east) -- (p.west |- a);
  \draw[arr] (b.east) -- (p.west |- b);
  \draw[arr] (c.east) -- (p.west |- c);
  \draw[arr] (d.east) -- (p.west |- d);
\end{tikzpicture}
\end{center}
\vspace{4pt}
Multiple candidates \textbf{compete} — best evidence wins.
\end{frame}

%% --- Slide 127: Stage 3 — Construction Loop ---
\begin{frame}
\frametitle{Stage 3 --- Construction Loop}
The 4-phase loop from Part III, applied per patch:
\vspace{8pt}
\begin{center}
\begin{tikzpicture}[
  phase/.style={draw=jugeoBlue, fill=jugeoLight!40, rounded corners,
                minimum width=2.2cm, minimum height=0.8cm,
                font=\footnotesize\bfseries, text=jugeoDark},
  arr/.style={-{Stealth[length=5pt]}, jugeoBlue, thick}
]
  \node[phase] (pr) at (0, 0) {PROPOSE};
  \node[phase] (no) at (3.5, 0) {NORMALIZE};
  \node[phase] (co) at (3.5,-1.8) {COMPARE};
  \node[phase] (se) at (0,-1.8) {SELECT};
  \draw[arr] (pr) -- (no);
  \draw[arr] (no) -- (co);
  \draw[arr] (co) -- (se);
  \draw[arr, jugeoOrange] (se) -- node[left, font=\scriptsize, text=jugeoOrange] {iterate} (pr);
\end{tikzpicture}
\end{center}
\vspace{6pt}
Picks the best candidate with \jugmark{full trust tracking}.
\end{frame}

%% --- Slide 128: Stage 4 — Descent and Glue ---
\begin{frame}
\frametitle{Stage 4 --- Descent and Glue}
\begin{block}{Overlap Check}
  Do all patches \textbf{agree} where they overlap?
\end{block}
\vspace{6pt}
\begin{center}
\begin{tikzpicture}[
  p/.style={draw=jugeoBlue, fill=jugeoLight!30, rounded corners,
            minimum width=2cm, minimum height=1cm, font=\footnotesize},
  ok/.style={draw=jugeoGreen!80, fill=jugeoGreen!15, rounded corners,
             font=\footnotesize\bfseries, inner sep=6pt},
  fail/.style={draw=jugeoRed!80, fill=jugeoRed!10, rounded corners,
               font=\footnotesize\bfseries, inner sep=6pt}
]
  \node[p] (a) at (0,0) {Patch A};
  \node[p] (b) at (3,0) {Patch B};
  \draw[jugeoGreen, very thick, <->] (a) -- node[above, font=\scriptsize, text=jugeoGreen]
    {agree} (b);
  \node[ok]   at (6.5, 0.5) {Glues $\Rightarrow$ global proof};
  \node[fail] at (6.5,-0.7) {Conflict $\Rightarrow$ obstruction};
\end{tikzpicture}
\end{center}
\vspace{4pt}
Extract \textcolor{jugeoRed}{obstructions} where patches disagree.\\
If everything glues, you have a \jugmark{global proof}.
\end{frame}

%% --- Slide 129: Stage 5 — Certificate Emit ---
\begin{frame}
\frametitle{Stage 5 --- Certificate Emit}
Annotate the winning solution with \jugmark{full provenance}:
\vspace{8pt}
\begin{center}
\renewcommand{\arraystretch}{1.4}
\begin{tabular}{l l}
  \toprule
  \textbf{Field} & \textbf{Content} \\
  \midrule
  Evidence     & Which solver / runtime data supported it \\
  Trust level  & \texttt{SOLVER\_DISCHARGED}, \texttt{RUNTIME}, \ldots \\
  Channels     & Source of each piece of evidence \\
  Coordinate   & Address in the semantic site \\
  \bottomrule
\end{tabular}
\end{center}
\vspace{6pt}
The certificate becomes part of the \textbf{Python file itself}.
\end{frame}

%% --- Slide 130: What Proof-Carrying Python Looks Like ---
\begin{frame}[fragile]
\frametitle{What Proof-Carrying Python Looks Like}
The scaffold starts with a \jugmark{coordinate function}:
\vspace{4pt}
\begin{block}{Site Address}
\begin{lstlisting}
def _spec_affine_program_0_coordinate():
    return "spec.affine.program.00"
\end{lstlisting}
\end{block}
\vspace{6pt}
\begin{itemize}
  \item The function's \textbf{address} in the semantic site
  \item Machine-readable, human-inspectable
  \item Links this file to the global verification map
\end{itemize}
\end{frame}

%% --- Slide 131: Scaffold — Support Function ---
\begin{frame}[fragile]
\frametitle{Scaffold --- Support Function}
The \jugmark{observable region} --- which inputs this function covers:
\vspace{4pt}
\begin{block}{Support}
\begin{lstlisting}
def _spec_affine_program_0_support(values):
    support = []
    for offset, value in enumerate(values):
        support.append((
            "spec.affine.program.00",
            offset,
            int(value)
        ))
    return tuple(support)
\end{lstlisting}
\end{block}
\vspace{4pt}
Returns a tuple of \texttt{(coordinate, offset, value)} triples.
\end{frame}

%% --- Slide 132: Scaffold — Descent Profile ---
\begin{frame}[fragile]
\frametitle{Scaffold --- Descent Profile}
\jugmark{Overlap metadata} --- how this function interacts with neighbors:
\vspace{4pt}
\begin{block}{Descent Profile}
\begin{lstlisting}
def _spec_affine_program_0_descent_profile(values):
    base = _spec_affine_program_0_coordinate()
    support = _spec_affine_program_0_support(values)
    return {
        "coordinate": base,
        "support":    support,
        "overlaps":   [],   # filled by glue stage
        "depth":      0,
    }
\end{lstlisting}
\end{block}
\vspace{4pt}
Used by Stage 4 (Descent and Glue) to check patch compatibility.
\end{frame}

%% --- Slide 133: The Implementation — Unchanged! ---
\begin{frame}[fragile]
\frametitle{The Implementation --- Unchanged!}
\begin{exampleblock}{Your Original Python}
\begin{lstlisting}
def solve(values, bias, mod, keep):
    total = 0
    for value in values:
        if int(value) % mod == keep:
            total += int(value) + bias
    return total
\end{lstlisting}
\end{exampleblock}
\vspace{6pt}
\begin{itemize}
  \item[\textcolor{jugeoGreen}{$\checkmark$}] No rewriting
  \item[\textcolor{jugeoGreen}{$\checkmark$}] No new syntax
  \item[\textcolor{jugeoGreen}{$\checkmark$}] Standard, importable, deployable
\end{itemize}
\end{frame}

%% --- Slide 134: The Spec — Also Python ---
\begin{frame}[fragile]
\frametitle{The Spec --- Also Python}
\begin{block}{Specification}
\begin{lstlisting}
def spec(result, values, bias, mod, keep):
    expected = sum(
        int(v) + bias
        for v in values
        if int(v) % mod == keep
    )
    return isinstance(result, int) and result == expected
\end{lstlisting}
\end{block}
\vspace{8pt}
\begin{center}
Both the \jugmark{code} and the \jugmark{spec} are plain Python.\\[4pt]
\textbf{Any developer can read them.}
\end{center}
\end{frame}

%% --- Slide 135: Declared Covers ---
\begin{frame}[fragile]
\frametitle{Declared Covers}
The \jugmark{geometric witness} --- inputs where the spec must hold:
\vspace{4pt}
\begin{block}{Cover Points}
\begin{lstlisting}
COVER = [
    {"args": [[-2, 0, 3, 6], -2, 3, 0]},
    {"args": [[0, 1, 2, 3, 4, 5], -2, 3, 0]},
    {"args": [[10, 20, 30], 0, 5, 0]},
    {"args": [[], 1, 2, 1]},
    {"args": [[7], 3, 7, 0]},
    {"args": [[0, 0, 0], 0, 1, 0]},
    {"args": [[1, -1, 1, -1], 10, 2, 1]},
    {"args": [[100], -100, 1, 0]},
    {"args": [[3, 6, 9, 12], 1, 3, 0]},
    {"args": [[2, 4, 6, 8], 0, 2, 0]},
]
\end{lstlisting}
\end{block}
\end{frame}

%% --- Slide 136: The Resulting Judgment ---
\begin{frame}
\frametitle{The Resulting Judgment}
\begin{center}
\begin{tikzpicture}[
  field/.style={draw=jugeoBlue!60, fill=jugeoLight!30, rounded corners,
                minimum width=10cm, minimum height=0.55cm,
                font=\scriptsize, text=jugeoDark, anchor=west},
  label/.style={font=\scriptsize\bfseries, text=jugeoBlue, anchor=east}
]
  \node[label] (l1) at (-0.15, 0.0) {coordinate};
  \node[field]  at ( 0.0, 0.0) {\texttt{"spec.affine.program.00"}};
  \node[label] (l2) at (-0.15,-0.75) {proposition};
  \node[field]  at ( 0.0,-0.75) {\texttt{solve} satisfies \texttt{spec} on COVER};
  \node[label] (l3) at (-0.15,-1.50) {evidence};
  \node[field, fill=jugeoGreen!15]  at ( 0.0,-1.50) {solver + runtime (10/10 cover points pass)};
  \node[label] (l4) at (-0.15,-2.25) {obligations};
  \node[field]  at ( 0.0,-2.25) {\texttt{[]} --- none remaining};
  \node[label] (l5) at (-0.15,-3.00) {obstructions};
  \node[field]  at ( 0.0,-3.00) {\texttt{[]} --- none found};
  \node[label] (l6) at (-0.15,-3.75) {trust};
  \node[field, fill=jugeoGreen!15]  at ( 0.0,-3.75) {\texttt{SOLVER\_DISCHARGED}};
  \node[label] (l7) at (-0.15,-4.50) {provenance};
  \node[field]  at ( 0.0,-4.50) {z3 solver $\times$ 10 inputs, runtime witness $\times$ 10 inputs};
\end{tikzpicture}
\end{center}
\end{frame}

%% --- Slide 137: When It Fails — Obstruction Example ---
\begin{frame}
\frametitle{When It Fails --- Obstruction Example}
Cover point 5 fails. What do we get?
\vspace{4pt}
\begin{center}
\begin{tikzpicture}[
  field/.style={draw=jugeoRed!50, fill=jugeoRed!8, rounded corners,
                minimum width=9.5cm, minimum height=0.55cm,
                font=\scriptsize, text=jugeoDark, anchor=west},
  label/.style={font=\scriptsize\bfseries, text=jugeoRed, anchor=east}
]
  \node[label] at (-0.15, 0.0) {coordinate};
  \node[field]  at ( 0.0, 0.0) {\texttt{"spec.affine.program.00"}};
  \node[label] at (-0.15,-0.75) {proposition};
  \node[field]  at ( 0.0,-0.75) {\texttt{spec(solve(*args), *args)} for cover point 5};
  \node[label] at (-0.15,-1.50) {actual};
  \node[field]  at ( 0.0,-1.50) {\texttt{solve([0,0,0], 0, 1, 0)} returned \texttt{1} (expected \texttt{0})};
  \node[label] at (-0.15,-2.25) {repair};
  \node[field]  at ( 0.0,-2.25) {Frontier: re-check \texttt{mod == 0} guard, \texttt{int()} cast};
\end{tikzpicture}
\end{center}
\vspace{6pt}
\begin{columns}[T]
\begin{column}{0.48\textwidth}
\begin{alertblock}{JuGeo}
  Structured \textbf{Obstruction} with\\
  coordinate, proposition, repair~frontier.
\end{alertblock}
\end{column}
\begin{column}{0.48\textwidth}
\begin{block}{F*}
  \texttt{\small(error) Unknown assertion failed.}\\[2pt]
  \textcolor{jugeoRed}{No repair guidance.}
\end{block}
\end{column}
\end{columns}
\end{frame}

%% --- Slide 138: The Four Proof Modes ---
\begin{frame}
\frametitle{The Four Proof Modes}
\begin{center}
\begin{tikzpicture}[
  mode/.style={draw=jugeoBlue, fill=jugeoLight!40, rounded corners=5pt,
               minimum width=4.5cm, minimum height=1cm,
               font=\footnotesize, text=jugeoDark, align=center},
  num/.style={circle, fill=jugeoBlue, text=white, font=\footnotesize\bfseries,
              inner sep=2pt, minimum size=0.5cm}
]
  \node[num]  (n1) at (-3.2, 1.2) {1};
  \node[mode] (m1) at ( 0.0, 1.2) {\textbf{Spec Satisfaction}\\Does program satisfy spec on cover?};
  \node[num]  (n2) at (-3.2, -0.2) {2};
  \node[mode] (m2) at ( 0.0,-0.2) {\textbf{Equivalence}\\Are two programs equal on cover?};
  \node[num]  (n3) at (-3.2, -1.6) {3};
  \node[mode] (m3) at ( 0.0,-1.6) {\textbf{Bug Detection}\\Does program have known bug patterns?};
  \node[num]  (n4) at (-3.2, -3.0) {4};
  \node[mode] (m4) at ( 0.0,-3.0) {\textbf{Relational Refinement}\\Compute a refined version};
\end{tikzpicture}
\end{center}
\end{frame}

%% --- Slide 139: Bug Detection — Six Classes ---
\begin{frame}
\frametitle{Bug Detection --- Six Classes}
Each detected as a \jugmark{structured obstruction}, not just a warning.
\vspace{6pt}
\begin{center}
\renewcommand{\arraystretch}{1.3}
\begin{tabular}{l l}
  \toprule
  \textbf{Bug Class} & \textbf{What It Catches} \\
  \midrule
  \texttt{bare-except}          & Catches \texttt{except:} without a type \\
  \texttt{identity-literal}     & Compares with \texttt{is} to a literal \\
  \texttt{late-binding-closure} & Closure captures loop variable by ref \\
  \texttt{mutable-default}      & Mutable default argument (\texttt{def f(x=[])}) \\
  \texttt{open-without-close}   & File handle never closed \\
  \texttt{shadow-builtin}       & Local name shadows a builtin \\
  \bottomrule
\end{tabular}
\end{center}
\vspace{6pt}
Every detection produces an \textcolor{jugeoRed}{Obstruction} with\\
coordinate, evidence, and \textbf{repair frontier}.
\end{frame}

%% --- Slide 140: The Proof Ships With The Code ---
\begin{frame}
\frametitle{The Proof Ships With The Code}
Any verifier can replay the proof from the Python file alone:
\vspace{8pt}
\begin{center}
\begin{tikzpicture}[
  vstep/.style={draw=jugeoBlue, fill=jugeoLight!40, rounded corners=4pt,
               minimum width=8cm, minimum height=0.7cm,
               font=\footnotesize, text=jugeoDark, align=left},
  num/.style={circle, fill=jugeoGreen!80, text=white,
              font=\footnotesize\bfseries, inner sep=2pt, minimum size=0.5cm},
  arr/.style={-{Stealth[length=4pt]}, jugeoBlue!60}
]
  \node[num]  (n1) at (-4.8, 0.0) {1};
  \node[vstep] (s1) at ( 0.0, 0.0) {Read \texttt{\_coordinate()} for the site address};
  \node[num]  (n2) at (-4.8,-1.1) {2};
  \node[vstep] (s2) at ( 0.0,-1.1) {Read \texttt{\_support()} for coverage};
  \node[num]  (n3) at (-4.8,-2.2) {3};
  \node[vstep] (s3) at ( 0.0,-2.2) {Re-run \texttt{spec} on every cover point};
  \node[num]  (n4) at (-4.8,-3.3) {4};
  \node[vstep] (s4) at ( 0.0,-3.3) {Inspect the judgment for trust \& provenance};
  \draw[arr] (s1) -- (s2);
  \draw[arr] (s2) -- (s3);
  \draw[arr] (s3) -- (s4);
\end{tikzpicture}
\end{center}
\vspace{8pt}
\begin{center}
\jugmark{No separate artifact.  No different language.}\\[4pt]
\textcolor{jugeoGreen}{\textbf{The proof is the code.}}
\end{center}
\end{frame}
